\documentclass[11pt, a4paper]{article}

% ── Packages ──────────────────────────────────────────────────────────────────
\usepackage[a4paper, top=2.5cm, bottom=2.5cm, left=2.5cm, right=2.5cm, headheight=28pt]{geometry}
\usepackage[utf8]{inputenc}
\usepackage[T1]{fontenc}
\usepackage{lmodern}
\usepackage{booktabs}
\usepackage{array}
\usepackage{tabularx}
\usepackage{longtable}
\usepackage{multirow}
\usepackage{colortbl}
\usepackage{xcolor}
\usepackage{titlesec}
\usepackage{enumitem}
\usepackage{hyperref}
\usepackage{fancyhdr}
\usepackage{parskip}
\usepackage{listings}
\usepackage[most]{tcolorbox}
\usepackage{fontawesome5}
\usepackage{microtype}
\usepackage{graphicx}
\usepackage{amsmath}

% ── Colors (KMUTNB brand) ──────────────────────────────────────────────────────
\definecolor{kmutnbBlue}{RGB}{0,56,116}
\definecolor{kmutnbGold}{RGB}{186,146,36}
\definecolor{lightGray}{RGB}{245,245,245}
\definecolor{midGray}{RGB}{200,200,200}
\definecolor{codeGray}{RGB}{248,248,248}
\definecolor{codeFrame}{RGB}{220,220,220}
\definecolor{tipteal}{RGB}{0,105,92}
\definecolor{warnAmber}{RGB}{121,85,0}
\definecolor{checkGreen}{RGB}{27,94,32}

% ── Section formatting ─────────────────────────────────────────────────────────
\titleformat{\section}
  {\large\bfseries\color{kmutnbBlue}}
  {\thesection.}{0.5em}{}
  [\color{kmutnbGold}\titlerule]

\titleformat{\subsection}
  {\normalsize\bfseries\color{kmutnbBlue}}
  {\thesubsection.}{0.5em}{}

\titleformat{\subsubsection}
  {\small\bfseries\color{kmutnbBlue}}
  {\thesubsubsection.}{0.5em}{}

\titlespacing*{\section}{0pt}{1.4ex plus .2ex}{0.6ex plus .1ex}
\titlespacing*{\subsection}{0pt}{1.0ex plus .2ex}{0.4ex plus .1ex}

% ── Header / Footer ────────────────────────────────────────────────────────────
\pagestyle{fancy}
\fancyhf{}
\fancyhead[L]{\small\color{kmutnbBlue}\textbf{KMUTNB -- Faculty of Engineering}}
\fancyhead[R]{\small\color{kmutnbBlue}Dept.\ of Electrical and Computer Engineering}
\fancyfoot[L]{\small\color{kmutnbBlue}Real-Time Operating Systems -- M.Eng.\ ECE}
\fancyfoot[C]{\small\thepage}
\fancyfoot[R]{\small\color{kmutnbBlue}Lab 01}
\renewcommand{\headrulewidth}{0.4pt}
\renewcommand{\headrule}{\hbox to\headwidth{%
  \color{kmutnbGold}\leaders\hrule height\headrulewidth\hfill}}
\renewcommand{\footrulewidth}{0.4pt}
\renewcommand{\footrule}{\hbox to\headwidth{%
  \color{midGray}\leaders\hrule height\footrulewidth\hfill}}

% ── Hyperlinks ─────────────────────────────────────────────────────────────────
\hypersetup{
  colorlinks = true,
  linkcolor  = kmutnbBlue,
  urlcolor   = kmutnbBlue,
  citecolor  = kmutnbBlue,
  pdftitle   = {Lab 01 — Toolchain and Environment Verification},
  pdfauthor  = {KMUTNB RTOS Course},
}

% ── Listings (C and bash) ──────────────────────────────────────────────────────
\lstdefinestyle{cstyle}{
  language        = C,
  basicstyle      = \ttfamily\small,
  keywordstyle    = \bfseries\color{kmutnbBlue},
  commentstyle    = \itshape\color{gray},
  stringstyle     = \color{tipteal},
  numberstyle     = \tiny\color{midGray},
  numbers         = left,
  numbersep       = 6pt,
  stepnumber      = 1,
  backgroundcolor = \color{codeGray},
  frame           = single,
  framerule       = 0.4pt,
  rulecolor       = \color{codeFrame},
  breaklines      = true,
  showstringspaces= false,
  tabsize         = 4,
  xleftmargin     = 1.5em,
  xrightmargin    = 0.5em,
  aboveskip       = 0.8em,
  belowskip       = 0.8em,
}

\lstdefinestyle{bashstyle}{
  language        = bash,
  basicstyle      = \ttfamily\small,
  keywordstyle    = \color{kmutnbBlue},
  commentstyle    = \itshape\color{gray},
  backgroundcolor = \color{black!5},
  frame           = single,
  framerule       = 0.4pt,
  rulecolor       = \color{codeFrame},
  breaklines      = true,
  showstringspaces= false,
  xleftmargin     = 1.5em,
  xrightmargin    = 0.5em,
  aboveskip       = 0.6em,
  belowskip       = 0.6em,
}

\lstdefinestyle{textstyle}{
  basicstyle      = \ttfamily\footnotesize,
  backgroundcolor = \color{black!4},
  frame           = single,
  framerule       = 0.4pt,
  rulecolor       = \color{codeFrame},
  breaklines      = true,
  xleftmargin     = 1.5em,
  xrightmargin    = 0.5em,
  aboveskip       = 0.6em,
  belowskip       = 0.6em,
}

% ── Callout boxes ──────────────────────────────────────────────────────────────
\tcbuselibrary{skins, breakable}

\newtcolorbox{tipbox}[1][]{
  enhanced, breakable,
  colback=tipteal!8, colframe=tipteal, coltitle=white,
  fonttitle=\bfseries\small, title={\faLightbulb\quad Tip#1},
  left=6pt, right=6pt, top=4pt, bottom=4pt,
  boxrule=0.8pt, arc=3pt,
}

\newtcolorbox{warnbox}[1][]{
  enhanced, breakable,
  colback=warnAmber!8, colframe=kmutnbGold, coltitle=white,
  fonttitle=\bfseries\small, title={\faExclamationTriangle\quad Note#1},
  left=6pt, right=6pt, top=4pt, bottom=4pt,
  boxrule=0.8pt, arc=3pt,
}

\newtcolorbox{checkpointbox}[1][]{
  enhanced,
  colback=kmutnbBlue!6, colframe=kmutnbBlue, coltitle=white,
  fonttitle=\bfseries\small, title={\faCheckCircle\quad Checkpoint#1},
  left=6pt, right=6pt, top=4pt, bottom=4pt,
  boxrule=0.8pt, arc=3pt,
}

\newtcolorbox{questionbox}[1][]{
  enhanced, breakable,
  colback=kmutnbGold!10, colframe=kmutnbGold, coltitle=white,
  fonttitle=\bfseries\small, title={\faQuestion\quad Question#1},
  left=6pt, right=6pt, top=4pt, bottom=4pt,
  boxrule=0.8pt, arc=3pt,
}

% ── Checkbox list ──────────────────────────────────────────────────────────────
\newlist{checklist}{itemize}{1}
\setlist[checklist]{label=\faSquare[regular], leftmargin=1.5em, noitemsep}

% ── Misc helpers ───────────────────────────────────────────────────────────────
\newcommand{\file}[1]{\texttt{\small #1}}
\newcommand{\api}[1]{\texttt{\small #1}}

% ══════════════════════════════════════════════════════════════════════════════
\begin{document}

% ── Title block ───────────────────────────────────────────────────────────────
\begin{center}
  {\color{kmutnbBlue}\rule{\linewidth}{2pt}}\\[6pt]
  {\LARGE\bfseries\color{kmutnbBlue} Laboratory 01}\\[4pt]
  {\large\bfseries Toolchain \& Environment Verification}\\[2pt]
  {\normalsize Real-Time Operating Systems \enspace\textbullet\enspace
   M.Eng.\ ECE \enspace\textbullet\enspace KMUTNB}\\[8pt]
  {\color{kmutnbGold}\rule{\linewidth}{1pt}}
\end{center}

\vspace{0.3cm}

\begin{tabular}{@{}p{3.8cm} p{10cm}@{}}
  \textbf{Platform}       & NXP FRDM-MCXN236 (ARM Cortex-M33 @ 150\,MHz PLL) \\[2pt]
  \textbf{RTOS}           & FreeRTOS via NXP MCUXpresso SDK (\texttt{freertos-kernel-upstream}) \\[2pt]
  \textbf{Toolchain}      & ARM GNU Toolchain 12 (\texttt{arm-none-eabi-gcc}) + Make \\[2pt]
  \textbf{Flash / Debug}  & LinkServer (\texttt{/Applications/LinkServer\_26.3.123/}) \\[2pt]
  \textbf{Estimated time} & 2--3 hours \\[2pt]
  \textbf{CLO mapping}    & CLO\,1 (foundations) \enspace CLO\,3 (FreeRTOS implementation) \\
\end{tabular}

\vspace{0.4cm}

% ── Objectives ────────────────────────────────────────────────────────────────
\section{Objectives}

By completing this lab you will be able to:

\begin{enumerate}[noitemsep, leftmargin=*]
  \item Build a FreeRTOS project using the NXP MCUXpresso SDK and ARM cross-compiler.
  \item Flash firmware and monitor the board using LinkServer.
  \item Verify that FreeRTOS creates and schedules tasks correctly via UART output
        and a physical LED.
  \item Observe how task priority affects scheduling --- the foundation for Module~2.
  \item Identify and prevent a race condition on a shared resource using a mutex.
\end{enumerate}

% ── Prerequisites ─────────────────────────────────────────────────────────────
\section{Prerequisites}

Before starting, confirm these are in place:

\begin{checklist}
  \item ARM GNU Toolchain 12 installed; \texttt{arm-none-eabi-gcc} available at\\
        \texttt{/opt/homebrew/Cellar/arm-gcc-bin@12/12.2.Rel1/bin}\\
        (or update \texttt{TOOLCHAIN} in \file{Makefile})
  \item NXP MCUXpresso SDK extracted; update \texttt{SDK\_ROOT} in \file{Makefile} if needed
  \item \textbf{LinkServer} installed: \texttt{/Applications/LinkServer\_26.3.123/LinkServer}
  \item FRDM-MCXN236 connected via the \textbf{MCU-LINK USB} port
  \item Python~3 + \texttt{pyserial} for serial monitoring (\texttt{pip3 install pyserial})
\end{checklist}

\begin{warnbox}
  This lab uses \textbf{LinkServer} (NXP's native flash/debug tool) ---
  not pyocd or J-Link. The board does \textbf{not} need its MCU-Link firmware
  reflashed. LinkServer works with the factory-default MCU-Link firmware.
\end{warnbox}

% ── Project Structure ─────────────────────────────────────────────────────────
\section{Project Structure}

\begin{lstlisting}[style=textstyle]
lab01_setup_verify/
  Makefile                      <- build, flash, debug targets
  src/
    main.c                      <- task definitions, LED init, scheduler start
    FreeRTOSConfig.h            <- kernel configuration
  <SDK_ROOT>/                   <- board support, FreeRTOS kernel, drivers (external)
    rtos/freertos/freertos-kernel-upstream/   <- FreeRTOS kernel
    devices/MCX/MCXN/MCXN236/               <- device startup + linker script
    examples/_boards/frdmmcxn236/            <- board.h, clock_config, pin_mux
\end{lstlisting}

The build pulls all SDK sources directly from \texttt{SDK\_ROOT} in the Makefile ---
there is nothing to clone or install separately.

\subsection{What the code does}

Two FreeRTOS tasks run continuously after the scheduler starts:

\renewcommand{\arraystretch}{1.3}
\begin{tabular}{@{} l c c p{6cm} @{}}
  \toprule
  \textbf{Task} & \textbf{Priority} & \textbf{Period} & \textbf{What it does} \\
  \midrule
  \texttt{vBlinkTask}     & 2 (HIGH) & 500\,ms  &
    Toggles the \textbf{physical red LED} (PIO4\_18) and prints \texttt{LED ON/OFF} \\
  \texttt{vHeartbeatTask} & 1 (LOW)  & 1000\,ms &
    Prints the FreeRTOS tick uptime in milliseconds \\
  \bottomrule
\end{tabular}
\renewcommand{\arraystretch}{1}

\medskip
A shared \textbf{UART mutex} (\api{xUartMutex}) prevents both tasks from calling
\api{PRINTF} simultaneously --- the simplest possible example of a critical section.

\begin{tipbox}
  \textbf{Hardware note:} The red LED is active-low on GPIO4 pin~18 (board pad N10).\\
  \texttt{LOGIC\_LED\_ON = 0} (GPIO low $\to$ LED on),
  \texttt{LOGIC\_LED\_OFF = 1} (GPIO high $\to$ LED off).\\
  \api{LED\_RED\_TOGGLE()} is a macro defined in the SDK's \file{board.h}.
\end{tipbox}

% ── Part A ─────────────────────────────────────────────────────────────────────
\section{Part A --- Build and Flash}

\subsection{Step 1 --- Build}

\begin{lstlisting}[style=bashstyle]
cd labs/lab01_setup_verify
make
\end{lstlisting}

Expected output (sizes will vary slightly):

\begin{lstlisting}[style=textstyle]
arm-none-eabi-gcc ... -o lab01.elf
arm-none-eabi-objcopy -O binary lab01.elf lab01.bin
arm-none-eabi-size lab01.elf
   text    data     bss     dec     hex filename
  31588      20   16060   47668    ba34 lab01.elf
\end{lstlisting}

\begin{warnbox}
  \textbf{Toolchain not found:} Edit the \texttt{TOOLCHAIN} variable at the top of
  \file{Makefile} to point to your \texttt{arm-none-eabi-gcc} installation directory.

  \medskip
  \textbf{SDK not found:} Edit \texttt{SDK\_ROOT} in \file{Makefile} to match
  where you extracted the NXP MCUXpresso SDK.
\end{warnbox}

\subsection{Step 2 --- Flash}

Connect the board via the MCU-LINK USB port, then:

\begin{lstlisting}[style=bashstyle]
make flash
\end{lstlisting}

This calls LinkServer to erase, program, and reset the MCU.
Look for \texttt{Finished writing Flash successfully} in the output.
The board starts running immediately after the reset.

\begin{tipbox}
  \textbf{Alternative --- pyocd:} pyocd is available in a virtualenv at\\
  \texttt{/Users/rus/Development/nxp\_test/.venv/bin/pyocd}\\
  Activate it first, then run: \texttt{pyocd flash lab01.elf -t mcxn236}
\end{tipbox}

% ── Part B ─────────────────────────────────────────────────────────────────────
\section{Part B --- Verify: UART Output and LED}

\subsection{Serial terminal}

The board exposes a USB-CDC virtual COM port over the MCU-LINK USB cable.

\textbf{Find the port:}
\begin{lstlisting}[style=bashstyle]
ls /dev/cu.usbmodem*
# Example output: /dev/cu.usbmodemLBAZMNNVFKEDG3
\end{lstlisting}

\textbf{Monitor with screen:}
\begin{lstlisting}[style=bashstyle]
screen /dev/cu.usbmodemXXXXXX 115200
# Press Ctrl-A then Ctrl-\ to quit
\end{lstlisting}

\textbf{Or with Python (clean, non-garbled capture):}
\begin{lstlisting}[style=bashstyle]
python3 -c "
import serial, time, sys
s = serial.Serial('/dev/cu.usbmodemXXXXXX', 115200, timeout=0.1)
s.reset_input_buffer()
end = time.time() + 10
while time.time() < end:
    d = s.read(256)
    if d: sys.stdout.buffer.write(d); sys.stdout.buffer.flush()
s.close()
"
\end{lstlisting}

\textbf{Windows:} PuTTY or Tera Term --- \texttt{COMx}, 115200\,baud, 8N1.

\subsection{Expected output}

\begin{lstlisting}[style=textstyle]
=== RTOS Lab 01: Toolchain & Environment Verify ===
SDK: BOARD_InitHardware + PRINTF  |  CPU: PLL 150 MHz
Tasks: Blink (500 ms) | Heartbeat (1000 ms)

[MAIN] Starting scheduler...

[BLINK] LED ON   (toggle #1)
[BLINK] LED OFF  (toggle #2)
[HB]    uptime = 1003 ms
[BLINK] LED ON   (toggle #3)
[BLINK] LED OFF  (toggle #4)
[HB]    uptime = 2007 ms
\end{lstlisting}

\begin{tipbox}
  \textbf{What to observe:}
  \begin{itemize}[noitemsep]
    \item \texttt{[BLINK]} prints \emph{twice} for every one \texttt{[HB]} line,
          reflecting 500\,ms vs.\ 1000\,ms periods.
    \item Uptime increments by approximately 1000\,ms per \texttt{[HB]} line.
    \item The \textbf{red LED} on the board blinks visibly at 2\,Hz
          (500\,ms on, 500\,ms off).
    \item Output never garbles --- the mutex serialises both tasks' \api{PRINTF} calls.
  \end{itemize}
\end{tipbox}

\begin{checkpointbox}
  Capture at least 5 seconds of terminal output \textbf{and} photograph the blinking LED.
\end{checkpointbox}

% ── Part C ─────────────────────────────────────────────────────────────────────
\section{Part C --- Experiments}

Work through these in order.  Record UART output for each one.

\subsection{Experiment 1 --- Priority swap}

In \file{main.c}, swap the task priorities:

\begin{lstlisting}[style=cstyle]
xTaskCreate(vBlinkTask,     "Blink", 512, NULL, 1, NULL); /* was 2 */
xTaskCreate(vHeartbeatTask, "HB",    512, NULL, 2, NULL); /* was 1 */
\end{lstlisting}

Rebuild (\texttt{make}), flash, and observe the serial output.

\begin{questionbox}
  Does the output order or timing change?  Why or why not?

  \textit{Hint: both tasks spend almost all their time blocked inside
  \api{vTaskDelay}.  Think about what happens when both unblock at the exact same tick.}
\end{questionbox}

\subsection{Experiment 2 --- Equal priorities}

Set both tasks to the same priority:

\begin{lstlisting}[style=cstyle]
xTaskCreate(vBlinkTask,     "Blink", 512, NULL, 1, NULL);
xTaskCreate(vHeartbeatTask, "HB",    512, NULL, 1, NULL);
\end{lstlisting}

\begin{questionbox}
  What scheduling policy does FreeRTOS apply when two tasks have equal priority and are
  both ready?  Where in \file{FreeRTOSConfig.h} is this policy enabled?
\end{questionbox}

\subsection{Experiment 3 --- Periods and CPU load}

Restore original priorities (Blink~= 2, HB~= 1).
Change \texttt{vBlinkTask}'s delay from 500\,ms to \textbf{50\,ms}:

\begin{lstlisting}[style=cstyle]
vTaskDelay(pdMS_TO_TICKS(50));
\end{lstlisting}

Rebuild and flash.  The red LED should now blink at 10\,Hz.

\begin{questionbox}
  Using the utilization formula $U = C/T$, estimate the CPU load of
  \texttt{vBlinkTask}. $C$~is the time spent executing (very short --- one
  \api{PRINTF} call); $T = 50$\,ms.

  Does the scheduler still meet the 1000\,ms heartbeat deadline?  Why?
\end{questionbox}

\subsection{Experiment 4 --- Remove the mutex (observe the race)}

Keep the 50\,ms Blink period from Experiment~3.  In both task bodies, comment out
the \api{xSemaphoreTake} / \api{xSemaphoreGive} lines around \api{PRINTF}:

\begin{lstlisting}[style=cstyle]
/* xSemaphoreTake(xUartMutex, portMAX_DELAY); */
PRINTF("[BLINK] LED %s  (toggle #%u)\r\n", state, (unsigned)count);
/* xSemaphoreGive(xUartMutex); */
\end{lstlisting}

\begin{questionbox}
  Do you observe garbled or interleaved output on the terminal?

  Why is this a \textbf{race condition} even on a \emph{single-core} MCU?
  What property of a mutex prevents it?
\end{questionbox}

\begin{warnbox}
  Restore the mutex calls before submitting.
\end{warnbox}

% ── Deliverables ───────────────────────────────────────────────────────────────
\section{Deliverables}

Submit the following to the course LMS by the posted deadline:

\renewcommand{\arraystretch}{1.3}
\begin{tabular}{@{} c p{4cm} p{8cm} @{}}
  \toprule
  \textbf{\#} & \textbf{Item} & \textbf{Details} \\
  \midrule
  1 & Terminal capture    & 5+ seconds of correct \texttt{[BLINK]} / \texttt{[HB]} output \\
  2 & LED photo           & Board with red LED visibly on or blinking \\
  3 & Answers to Q1--Q4  & A few sentences each; reference the code or output \\
  4 & Modified \file{main.c} & From Experiment~3 (50\,ms Blink, mutex restored) \\
  \bottomrule
\end{tabular}
\renewcommand{\arraystretch}{1}

% ── Key Concepts ───────────────────────────────────────────────────────────────
\section{Key Concepts Introduced}

\renewcommand{\arraystretch}{1.3}
\begin{tabular}{@{} p{6.5cm} p{7cm} @{}}
  \toprule
  \textbf{Concept} & \textbf{Where you saw it} \\
  \midrule
  Preemptive, priority-based scheduling    & Priority swap (Experiment~1) \\
  Time-slicing among equal-priority tasks  & Equal priority (Experiment~2) \\
  CPU utilization $U = C/T$               & Short-period Blink (Experiment~3) \\
  Race condition on a shared resource      & Mutex removal (Experiment~4) \\
  Mutex as a critical-section guard        & \api{xUartMutex} throughout \\
  Periodic task with \api{vTaskDelay}      & \texttt{vBlinkTask}, \texttt{vHeartbeatTask} \\
  \bottomrule
\end{tabular}
\renewcommand{\arraystretch}{1}

\medskip
These are the building blocks for Module~2 (scheduling theory) and Module~3
(synchronisation).

% ── Implementation Notes ───────────────────────────────────────────────────────
\section{Implementation Notes}

\subsection{Why \texttt{configRUN\_FREERTOS\_SECURE\_ONLY 1}}

The MCXN236 boots and runs entirely in the ARM \textbf{secure world}.
The FreeRTOS ARM\_CM33\_NTZ port must know this to set the correct
\texttt{EXC\_RETURN} value when launching the first task:

\renewcommand{\arraystretch}{1.3}
\begin{tabular}{@{} l l c p{5cm} @{}}
  \toprule
  \textbf{Config} & \textbf{EXC\_RETURN} & \textbf{Bit[6]} & \textbf{Result} \\
  \midrule
  \texttt{configRUN\_FREERTOS\_SECURE\_ONLY 0} (default)
    & \texttt{0xFFFFFFBC} & 0 -- non-secure
    & Immediate fault: CPU is actually in secure mode \\
  \texttt{configRUN\_FREERTOS\_SECURE\_ONLY 1}
    & \texttt{0xFFFFFFFD} & 1 -- secure
    & First task starts correctly \\
  \bottomrule
\end{tabular}
\renewcommand{\arraystretch}{1}

\medskip
Without this setting the scheduler appears to start (the banner prints) but
no tasks ever execute --- the symptom is total silence after
\texttt{[MAIN] Starting scheduler...}

\subsection{Why the SDK FreeRTOS kernel}

This lab uses the NXP SDK's \texttt{freertos-kernel-upstream}
(inside \texttt{SDK\_ROOT/rtos/freertos/}) rather than a standalone clone.
The SDK kernel is the version NXP has tested against this chip and BSP.
A development-branch clone may have incompatible \file{portasm.c} changes that
interact badly with the secure-world boot sequence.

\subsection{LED hardware}

\renewcommand{\arraystretch}{1.3}
\begin{tabular}{@{} l c c c c @{}}
  \toprule
  \textbf{LED} & \textbf{GPIO} & \textbf{Pin} & \textbf{Board pad} & \textbf{Active level} \\
  \midrule
  Red   & GPIO4 & 18 & N10 & Low (0 = ON) \\
  Blue  & GPIO4 & 17 & R9  & Low (0 = ON) \\
  Green & GPIO4 & 19 & R10 & Low (0 = ON) \\
  \bottomrule
\end{tabular}
\renewcommand{\arraystretch}{1}

\medskip
All three share GPIO4. The pin mux is set to \texttt{kPORT\_MuxAlt0}
(GPIO function) in \file{main.c} before the scheduler starts,
using \api{PORT\_SetPinConfig} and \api{GPIO\_PinInit}.

% ── Reference ─────────────────────────────────────────────────────────────────
\section{Reference}

\begin{tabular}{@{} p{6.5cm} p{7cm} @{}}
  \toprule
  \textbf{Resource} & \textbf{Relevant sections} \\
  \midrule
  Barry, \textit{Mastering the FreeRTOS Kernel}
    & Ch.~1 (tasks), Ch.~3 (queue/semaphore basics) \\
  Laplante, \textit{Real-Time Systems Design}
    & Ch.~1--2 (task model, deadline classes) \\
  FreeRTOS documentation
    & \api{xTaskCreate}, \api{vTaskDelay}, \api{xSemaphoreCreateMutex} \\
  NXP AN13036
    & FreeRTOS on Cortex-M33 with TrustZone \\
  SDK \file{board.h}
    & \api{LED\_RED\_TOGGLE}, \api{LOGIC\_LED\_ON/OFF} macros \\
  \bottomrule
\end{tabular}

\end{document}
